%! Matrices as Functions — Unit 4, Lesson 4.13 — Warm-Up (Answer Key)
\documentclass[10pt]{article}
\usepackage{precalculus-article}
\usepackage{precalculus-key}   % pulls in -boxes; do NOT also load -boxes
\newcommand{\TallMath}[1]{$\displaystyle #1\rule[-1.4em]{0pt}{3.2em}$}

\begin{document}

\pageheader{Unit 4, Lesson 4.13}{Warm-Up --- Answer Key}
\namedateperiod

\vspace{0.12in}

\begin{notesbox}{Spiral Review --- Key}
\small

\textbf{1.}\quad Compute each matrix--vector product.

\vspace{0.06in}
\hspace{1em}
(a)\quad
$\begin{bmatrix}0 & -1 \\ 1 & 0\end{bmatrix}
\begin{bmatrix}1 \\ 0\end{bmatrix}
= $ \ans{$\begin{bmatrix}0 \\ 1\end{bmatrix}$}
\qquad\qquad
(b)\quad
$\begin{bmatrix}0 & -1 \\ 1 & 0\end{bmatrix}
\begin{bmatrix}0 \\ 1\end{bmatrix}
= $ \ans{$\begin{bmatrix}-1 \\ 0\end{bmatrix}$}

\vspace{0.2in}

\textbf{2.}\quad Find the determinant of $A = \begin{bmatrix}2 & 0 \\ 0 & 3\end{bmatrix}$.
Then explain what $|\det A|$ tells you about the area of any region transformed by $A$.

\vspace{0.06in}
\hspace{1em} $\det A = $ \ans{$6$}

\vspace{0.06in}
\hspace{1em} Area interpretation: \ans{The transformation scales all areas by a factor of $6$.}

\vspace{0.2in}

\textbf{3.}\quad Let $A = \begin{bmatrix}1 & 2 \\ 0 & 1\end{bmatrix}$.
Compute $A^{-1}$ using the $2\times2$ inverse formula
$\begin{bmatrix}a&b\\c&d\end{bmatrix}^{-1}=\dfrac{1}{ad-bc}\begin{bmatrix}d&-b\\-c&a\end{bmatrix}$.

\vspace{0.1in}
\hspace{1em} $A^{-1} = $ \ans{$\begin{bmatrix}1 & -2 \\ 0 & 1\end{bmatrix}$}

\vspace{0.06in}
\hspace{1em} Verify: \ans{$\begin{bmatrix}1&2\\0&1\end{bmatrix}\begin{bmatrix}1&-2\\0&1\end{bmatrix}=\begin{bmatrix}1&0\\0&1\end{bmatrix}=I$. \checkmark}

\end{notesbox}

\begin{teachernote}
Q1 previews the main idea: the columns of $\bigl[\begin{smallmatrix}0&-1\\1&0\end{smallmatrix}\bigr]$
are exactly the images of $\hat{\imath}$ and $\hat{\jmath}$ under the $90°$ rotation --- a direct
bridge to Step 1 of the lesson. If students struggle with Q3, remind them $\det A = 1\cdot1 - 2\cdot0 = 1$,
so the inverse formula simplifies to just swapping the diagonal and negating the off-diagonal.
\end{teachernote}

\end{document}
